\appendix
\chapter{Codes and Automation Scripts}
\label{ch:appendix_codes}

\section{DoE Case Generator Script (\texttt{doe\_generator.py})}
\label{sec:app_doe_gen}
\begin{lstlisting}[language=Python, caption={Python script for automated generation of the 24 OpenFOAM DoE cases.}]
import os
import shutil
import subprocess
import itertools

# ==================== 1. DoE PARAMETERS ====================

# Kinematic viscosity (nu) [m^2/s] at 20 C, 1 atm (NIST REFPROP)
fluids = {
    "H2": "1.05e-04",
    "N2": "1.51e-05"
}

# Velocity @ inlet [m/s]
fluxes = ["15", "40", "80"]

# Obturators
geometries = ["3Fori_1p5mm.stl", "3Fori_3mm.stl", "6Fori_1p5mm.stl", "6Fori_3mm.stl"]

# Folder definition
base_case_dir = r"C:\Users\andre\Desktop\Tesi\Simulazioni\CASE_BASE"
doe_run_dir = r"C:\Users\andre\Desktop\Tesi\Simulazioni"
geom_source_dir = r"C:\Users\andre\Desktop\Tesi\Simulazioni\Geometrie_STL"

# ==================== 2. HELPER FUNCTION ====================

def replace_in_file(filepath, search_text, replace_text):
    with open(filepath, 'r') as file:
        file_data = file.read()
    file_data = file_data.replace(search_text, replace_text)
    with open(filepath, 'w') as file:
        file.write(file_data)

# ==================== 3. GENERATE AND SUBMIT CASES ====================

if not os.path.exists(doe_run_dir):
    os.makedirs(doe_run_dir)

for fluid, nu_val in fluids.items():
    for flux in fluxes:
        for geom in geometries:
            geom_name = geom.replace('.stl', '')
            case_name = f"Sim_{fluid}_Flux{flux}_{geom_name}"
            case_path = os.path.join(doe_run_dir, case_name)

            print(f"Generating case: {case_name}")

            if os.path.exists(case_path):
                shutil.rmtree(case_path)
            shutil.copytree(base_case_dir, case_path)

            src_stl = os.path.join(geom_source_dir, geom)
            dst_stl = os.path.join(case_path, "constant", "geometry", "wallObt.stl")
            shutil.copyfile(src_stl, dst_stl)

            init_cond_path = os.path.join(case_path, "0", "include", "initialConditions")
            replace_in_file(init_cond_path, "FLUX_VALUE", flux)

            transport_path = os.path.join(case_path, "constant", "physicalProperties")
            replace_in_file(transport_path, "NU_VALUE", nu_val)

            # subprocess.run(["sbatch", "run_aries_2.sh"], cwd=case_path)

print("DoE Generation Complete. 24 cases prepared.")
\end{lstlisting}

\section{SLURM Execution and Maintenance Scripts (Bash)}
\label{sec:app_slurm_scripts}
\begin{lstlisting}[language=bash, caption={Full meshing and solving SLURM submission script (\texttt{run\_aries\_2.sh}).}]
#!/bin/bash
#SBATCH --account=fluids
#SBATCH --partition=high
#SBATCH --time=23:59:59
#SBATCH --nodes=1
#SBATCH --ntasks-per-node=52
#SBATCH --mem=490g
#SBATCH --job-name=CEDEM
#SBATCH --exclusive
#SBATCH --exclude=cnode[06,17,19]

module load gcc-12.1.1/gcc
module load gcc-12.1.1/ompi-4.1.5_nccl

source $HOME/OpenFOAM-10/etc/bashrc

TOTAL_TASKS=$((SLURM_NTASKS_PER_NODE * SLURM_NNODES))
echo "numberOfSubdomains $TOTAL_TASKS;" > system/numProcs

blockMesh
surfaceFeatures
decomposePar -copyZero -force
mpirun -np $TOTAL_TASKS snappyHexMesh -overwrite -parallel > out_mesh
mpirun -np $TOTAL_TASKS simpleFoam -parallel > out_run
\end{lstlisting} %

\section{Remote Post-Processing Extraction (\texttt{download\_postprocessing.py})}
\label{sec:app_scp_script}
\begin{lstlisting}[language=Python, caption={Python script utilizing \texttt{scp} to selectively download \texttt{postProcessing} folders from ARIES.}]
import os

hpc_user = "fluids11"
hpc_host = "hpc.unimo.lan"
remote_base_dir = "/unimore_home/fluids11/SIM_DoE_1610"

local_base_dir = os.path.dirname(os.path.abspath(__file__))

sim_folders = [f for f in os.listdir(local_base_dir) 
               if os.path.isdir(os.path.join(local_base_dir, f)) and f.startswith("Sim_")]

for case in sim_folders:
    remote_path = f"{hpc_user}@{hpc_host}:{remote_base_dir}/{case}/postProcessing"
    local_case_dir = os.path.join(local_base_dir, case)
    print(f"\n---> Fetching data for: {case}")
    scp_command = f'scp -r {remote_path} "{local_case_dir}"'
    os.system(scp_command)
\end{lstlisting} %

\section{Batch Plotly Report and Master Dashboard Generator (\texttt{generate\_reports.py})}
\label{sec:app_plotly_script}
\begin{lstlisting}[language=Python, caption={Automated data parser, Plotly report generator, and Master HTML Dashboard builder.}]
import os
import re
import plotly.graph_objects as go

def parse_scalar_data(filepath):
    times, values = [], []
    with open(filepath, 'r') as f:
        for line in f:
            if line.startswith('#') or not line.strip():
                continue
            parts = line.split()
            if len(parts) >= 2:
                times.append(float(parts[0]))
                values.append(float(parts[1]))
    return times, values

def parse_vector_data(filepath):
    times, x, y, z = [], [], [], []
    with open(filepath, 'r') as f:
        for line in f:
            if line.startswith('#') or not line.strip():
                continue
            parts = line.replace('(', '').replace(')', '').split()
            if len(parts) >= 4:
                times.append(float(parts[0]))
                x.append(float(parts[1]))
                y.append(float(parts[2]))
                z.append(float(parts[3]))
    return times, x, y, z

def parse_yplus(filepath):
    data = {}
    with open(filepath, 'r') as f:
        for line in f:
            if line.startswith('#') or not line.strip():
                continue
            parts = line.split()
            if len(parts) >= 5:
                t = float(parts[0])
                patch = parts[1]
                avg_yplus = float(parts[4])
                if patch not in data:
                    data[patch] = {'time': [], 'avg': []}
                data[patch]['time'].append(t)
                data[patch]['avg'].append(avg_yplus)
    return data

def parse_forces(filepath):
    times, fx, fy, fz = [], [], [], []
    regex = re.compile(r"([0-9\.eE+-]+)\s+\(\(([^\)]+)\)\s+\(([^\)]+)\)\)")
    with open(filepath, 'r') as f:
        for line in f:
            if line.startswith('#') or not line.strip():
                continue
            match = regex.search(line)
            if match:
                t = float(match.group(1))
                p_force = [float(v) for v in match.group(2).split()]
                v_force = [float(v) for v in match.group(3).split()]
                times.append(t)
                fx.append(p_force[0] + v_force[0])
                fy.append(p_force[1] + v_force[1])
                fz.append(p_force[2] + v_force[2])
    return times, fx, fy, fz

local_base_dir = r"C:\Users\andre\Desktop\Tesi\Simulazioni"
sim_folders = [f for f in os.listdir(local_base_dir) 
               if os.path.isdir(os.path.join(local_base_dir, f)) and f.startswith("Sim_")]
dashboard_links = []

for case in sorted(sim_folders):
    case_path = os.path.join(local_base_dir, case)
    pp_dir = os.path.join(case_path, "postProcessing")
    if not os.path.exists(pp_dir):
        continue

    file_in_p   = os.path.join(pp_dir, "inlet_pressure_avg", "0", "surfaceFieldValue.dat")
    file_out_p  = os.path.join(pp_dir, "outlet_pressure_avg", "0", "surfaceFieldValue.dat")
    file_flux   = os.path.join(pp_dir, "inlet_volume_flux", "0", "surfaceFieldValue.dat")
    file_yplus  = os.path.join(pp_dir, "yPlus1", "0", "yPlus.dat")
    file_forces = os.path.join(pp_dir, "forceCoeffs_Obt", "0", "forces.dat")
    
    t_in_p, in_p = parse_scalar_data(file_in_p)
    t_out_p, out_p = parse_scalar_data(file_out_p)
    t_flux, f_x, f_y, f_z = parse_vector_data(file_flux)
    yplus_data = parse_yplus(file_yplus)
    t_forces, fx, fy, fz = parse_forces(file_forces)
    # Figures generated and saved to CFD_Report_<case>.html (see Section 4.2)
\end{lstlisting} %